% GENERATED FILE. Edit canonical lessons, then run npm run generate:books.
% canonical-source-hash: fnv1a64:084c61e2ffda1045
% canonical-lessons: ZH-C137-jieri, ZH-W137-jie-joint, ZH-C137-chunjie, ZH-C137-guonian, ZH-C137-fangjia, ZH-C137-jiaqi

\chapter{Holidays}
\label{ch:zh-holidays}
\hlchaptermodality{137}

\begin{chapteropening}
\textbf{By the end of this chapter:} \emph{I can talk about festivals, the Spring Festival, the New Year and time off.}

Complete the last lesson of chapter 137: I can talk about festivals, the Spring Festival, the New Year and time off.
\end{chapteropening}

\section[\texorpdfstring{\zh{节日} (jiérì)}{节日 (jiérì)}]{\zh{节日} (jiérì) — festival}

\label{lesson:ZH-C137-jieri}



\begin{quote}
Before the new one: say the Mandarin for to thunder, then the Mandarin for pleasantly cool.
\end{quote}

\subsection*{You'll want to know: \zh{节日}}
\textbf{\zh{节日}} — \emph{jiérì} — \textquotedblleft{}festival\textquotedblright{}.

\begin{quote}
\textbf{\zh{今天是什么节日}} \emph{jīntiān shì shénme jiérì} — What festival is it today?
\end{quote}

The new character, in pieces: \zh{节} — the grass piece on top, a bar crossed by two short verticals, and under it a bar that turns down into a hook, with a long vertical through it. The next lesson puts it in your hand.

\begin{cousinweb}
\textbf{\zh{节日}} (\emph{jiérì}) is built from \zh{节} 'a joint' and \zh{日} 'day', both inherited from Middle Chinese, \emph{tset} and \emph{nyit} (Guangyun, in Baxter's transcription); the dictionary quotes an older book for the first: the joint of a bamboo. A festival as a joint in the year is the likely image, though it is not given there.
\end{cousinweb}

\begin{grammarlens}[title={What you've built}]
One more word for times of day.
\end{grammarlens}

\subsection*{Your turn}
\begin{itemize}
\raggedright
  \item \emph{Say it:} \emph{jiérì}
  \item \emph{Say it:} \emph{jiérì}, once more
  \item \emph{Say it:} \emph{jīntiān shì shénme jiérì} — What festival is it today?
  \item \emph{Recall:} say the Mandarin for to thunder, then the Mandarin for pleasantly cool, then say \emph{jiérì} again
\end{itemize}

\begin{tcolorbox}[breakable,colback=teal!4,colframe=teal!35!black,title={Before you move on}]
What does \zh{节日} mean? (\textquotedblleft{}Festival\textquotedblright{}.) Say it once more.
\end{tcolorbox}
\section[\texorpdfstring{\zh{节} (jié)}{节 (jié)}]{\zh{节} — the grass piece over a hook}

\label{lesson:ZH-W137-jie-joint}



\begin{quote}
The word first, then the character inside it.

\begin{itemize}
\raggedright
  \item \emph{Recall:} say \emph{jiérì} — \textbf{R1}, one lesson back, before it decays
\end{itemize}

\begin{quote}
\zh{菜} \zh{日}
\end{quote}

\emph{Read it:} each one aloud, with nothing said first — \emph{cài}, \emph{rì}
\end{quote}

\subsection*{Script — the grass piece over a hook}
\hlblockfigure{\detokenize{figures/ZH-W137-jie-joint-filmstrip.pdf}}{How \zh{节} is written, stroke by stroke}

\begin{quote}
\zh{节}
\end{quote}

\textbf{Two pieces, top to bottom.}

\textbf{The grass piece, three strokes}, as on top of \zh{菜}: a horizontal across, then a short vertical on the left and a short vertical on the right, both crossing it.

\textbf{The lower piece, two strokes.} A horizontal that turns down and ends in a hook to the left, then a long vertical down through it to the foot.

3 + 2 = five.

\subsection*{Your turn}
\begin{itemize}
\raggedright
  \item \emph{Write it:} \zh{节}, all five, with the model in front of you
  \item \emph{Cover:} write it again from memory, and count the pieces, not the strokes
\end{itemize}

\begin{tcolorbox}[breakable,colback=teal!4,colframe=teal!35!black,title={Before you move on}]
How many pieces in \zh{节}? (\textbf{Two}.) How many strokes in the grass piece? (\textbf{Three}.) Which stroke is the last? (\textbf{The long vertical}.)

Source: \href{https://raw.githubusercontent.com/chanind/hanzi-writer-data/68d10a4b21150cae5e1ebbd223eed289cf32d90c/data/\%E8\%8A\%82.json}{Hanzi Writer Data \zh{节}, pinned stroke paths and medians}.
\end{tcolorbox}
\section[\texorpdfstring{\zh{春节} (Chūnjié)}{春节 (Chūnjié)}]{\zh{春节} (Chūnjié) — the Spring Festival; Chinese New Year}

\label{lesson:ZH-C137-chunjie}



\begin{quote}
Before the new one: say the Mandarin for pleasantly cool, then the Mandarin for festival.
\end{quote}

\subsection*{You'll want to know: \zh{春节}}
\textbf{\zh{春节}} — \emph{Chūnjié} — \textquotedblleft{}the Spring Festival; Chinese New Year\textquotedblright{}.

\begin{quote}
\textbf{\zh{春节前大家都打扫}} \emph{Chūnjié qián dàjiā dōu dǎsǎo} — Before the festival, everybody cleans the house.
\end{quote}

\begin{cousinweb}
\textbf{\zh{春节}} (\emph{Chūnjié}) is built from \zh{春} 'spring' and \zh{节} 'festival', both inherited from Middle Chinese, \emph{tsyhwin} and \emph{tset} (Guangyun, in Baxter's transcription): the festival at the head of spring.
\end{cousinweb}

\begin{grammarlens}[title={What you've built}]
One more word for times of day.
\end{grammarlens}

\subsection*{Your turn}
\begin{itemize}
\raggedright
  \item \emph{Say it:} \emph{Chūnjié}
  \item \emph{Say it:} \emph{Chūnjié}, once more
  \item \emph{Say it:} \emph{Chūnjié qián dàjiā dōu dǎsǎo} — Before the festival, everybody cleans the house.
  \item \emph{Recall:} say the Mandarin for pleasantly cool, then the Mandarin for festival, then say \emph{Chūnjié} again
\end{itemize}

\begin{tcolorbox}[breakable,colback=teal!4,colframe=teal!35!black,title={Before you move on}]
What does \zh{春节} mean? (\textquotedblleft{}The Spring Festival; Chinese New Year\textquotedblright{}.) Say it once more.
\end{tcolorbox}
\section[guònián]{guònián — to celebrate the New Year}

\label{lesson:ZH-C137-guonian}



\begin{quote}
Before the new one: say the Mandarin for festival, then the Mandarin for Chinese New Year.
\end{quote}

\subsection*{You'll want to know: guònián}
\emph{guònián} — \textquotedblleft{}to celebrate the New Year\textquotedblright{}.

\begin{quote}
\emph{guònián wǒmen chī jiǎozi} — We eat dumplings to see in the new year.
\end{quote}

Its characters come later; for now it is a word for the ear.

\begin{cousinweb}
\emph{guònián} is built from \emph{guò} 'to pass' and \zh{年} 'year', both inherited from Middle Chinese, \emph{kwaH} and \emph{nen} (Guangyun, in Baxter's transcription). The pair is not given there; seeing the year through is the plain picture, but how it came to mean celebrating is not settled.
\end{cousinweb}

\begin{grammarlens}[title={What you've built}]
One more word for times of day.
\end{grammarlens}

\subsection*{Your turn}
\begin{itemize}
\raggedright
  \item \emph{Say it:} \emph{guònián}
  \item \emph{Say it:} \emph{guònián}, once more
  \item \emph{Say it:} \emph{guònián wǒmen chī jiǎozi} — We eat dumplings to see in the new year.
  \item \emph{Recall:} say the Mandarin for festival, then the Mandarin for Chinese New Year, then say \emph{guònián} again
\end{itemize}

\begin{tcolorbox}[breakable,colback=teal!4,colframe=teal!35!black,title={Before you move on}]
What does guònián mean? (\textquotedblleft{}To celebrate the New Year\textquotedblright{}.) Say it once more.
\end{tcolorbox}
\section[fàngjià]{fàngjià — to have time off; to go on holiday}

\label{lesson:ZH-C137-fangjia}



\begin{quote}
Before the new one: say the Mandarin for Chinese New Year, then the Mandarin for to celebrate the New Year.
\end{quote}

\subsection*{You'll want to know: fàngjià}
\emph{fàngjià} — \textquotedblleft{}to have time off; to go on holiday\textquotedblright{}.

\begin{quote}
\emph{míngtiān wǒmen fàngjià} — We're off tomorrow.
\end{quote}

Its characters come later; for now it is a word for the ear.

\begin{cousinweb}
\emph{fàngjià} is built from \emph{fàng} 'to let go' and \emph{jià} 'leave', both inherited from Middle Chinese, \emph{pjangH} and \emph{kaeH} (Guangyun, in Baxter's transcription), which glosses the second 'to borrow' and 'leave from duty'.
\end{cousinweb}

\begin{grammarlens}[title={What you've built}]
One more word for times of day.
\end{grammarlens}

\subsection*{Your turn}
\begin{itemize}
\raggedright
  \item \emph{Say it:} \emph{fàngjià}
  \item \emph{Say it:} \emph{fàngjià}, once more
  \item \emph{Say it:} \emph{míngtiān wǒmen fàngjià} — We're off tomorrow.
  \item \emph{Recall:} say the Mandarin for Chinese New Year, then the Mandarin for to celebrate the New Year, then say \emph{fàngjià} again
\end{itemize}

\begin{tcolorbox}[breakable,colback=teal!4,colframe=teal!35!black,title={Before you move on}]
What does fàngjià mean? (\textquotedblleft{}To have time off; to go on holiday\textquotedblright{}.) Say it once more.
\end{tcolorbox}
\section[jiàqī]{jiàqī — holidays; time off}

\label{lesson:ZH-C137-jiaqi}



\begin{quote}
Before the new one: say the Mandarin for to celebrate the New Year, then the Mandarin for to have time off.
\end{quote}

\subsection*{You'll want to know: jiàqī}
\emph{jiàqī} — \textquotedblleft{}holidays; time off\textquotedblright{}.

\begin{quote}
\emph{fàngjià le jiàqī nǐ xiǎng qù nǎr} — We're on holiday now; where would you like to go?
\end{quote}

Its characters come later; for now it is a word for the ear.

\begin{cousinweb}
\emph{jiàqī} is built from \emph{jià} 'leave' and \zh{期} 'a set time', both inherited from Middle Chinese, \emph{kaeH} and \emph{gi} (Guangyun, in Baxter's transcription).
\end{cousinweb}

\begin{grammarlens}[title={What you've built}]
One more word for times of day.
\end{grammarlens}

\subsection*{Your turn}
\begin{itemize}
\raggedright
  \item \emph{Say it:} \emph{jiàqī}
  \item \emph{Say it:} \emph{jiàqī}, once more
  \item \emph{Say it:} \emph{fàngjià le jiàqī nǐ xiǎng qù nǎr} — We're on holiday now; where would you like to go?
  \item \emph{Recall:} say the Mandarin for to celebrate the New Year, then the Mandarin for to have time off, then say \emph{jiàqī} again
\end{itemize}

\begin{tcolorbox}[breakable,colback=teal!4,colframe=teal!35!black,title={Before you move on}]
What does jiàqī mean? (\textquotedblleft{}Holidays; time off\textquotedblright{}.) Say it once more.
\end{tcolorbox}
